\documentclass[11pt]{article}

\usepackage{amsmath}
\usepackage{amssymb}
\usepackage{amsthm}
\usepackage{url}
\usepackage[margin=1.15in]{geometry}
\usepackage{booktabs}
\usepackage{array}
\usepackage[colorlinks=true, linkcolor=blue, citecolor=blue, urlcolor=blue]{hyperref}

\newtheorem{theorem}{Theorem}[section]
\newtheorem{lemma}[theorem]{Lemma}
\newtheorem{proposition}[theorem]{Proposition}
\newtheorem{corollary}[theorem]{Corollary}

\theoremstyle{definition}
\newtheorem{definition}[theorem]{Definition}

\theoremstyle{remark}
\newtheorem{remark}[theorem]{Remark}
\newtheorem*{remark*}{Remark}

\emergencystretch=1.5em

\newcommand{\Num}{\operatorname{Num}}
\newcommand{\NS}{\operatorname{NS}}
\newcommand{\End}{\operatorname{End}}
\newcommand{\Fq}{\mathbb{F}_q}
\newcommand{\OX}{\mathcal{O}_X}

\title{The Castelnuovo--Severi/Hodge Index Inequality on a Surface from Riemann--Roch and Ampleness: A Non-Circularity Re-derivation}
\author{Kunal Tyagi\thanks{%
\texttt{hello.jay.tyagi@gmail.com}. \textbf{Use of AI.} The mathematics in this
paper --- the derivations, the computations, the verification suite, and the
text --- was produced by Claude (Anthropic) working under the author's direction
inside a structured research program; the author set the program's objectives,
adjudicated its decisions, and is responsible for the content. Claude is a tool and
is not an author. The complete program record, including every session log, review
and correction, is public: \url{https://github.com/x67ai/riemann-rh-program}.}}
\date{}

\begin{document}

\maketitle

\begin{abstract}
This note is a re-derivation of standard material for the record; zero novelty is claimed.
Every statement is graduate-textbook algebraic geometry. Its value is the audit trail: the
Hodge-index/Castelnuovo--Severi inequality on a smooth projective surface --- the positivity
engine of the only \emph{positivity mechanism} that has ever proved a Riemann hypothesis,
namely Weil's proof of RH for curves over finite fields --- is proved here in full from
Riemann--Roch on the surface, Serre duality, and ampleness, with no zeta-function,
$L$-function, or RH input anywhere. The proof is valid over an algebraically closed field of
arbitrary characteristic. We itemize the provenance of every input in a step-by-step ledger,
verifying that no zeta function, point count, explicit-formula positivity, or statement
equivalent to any Riemann hypothesis appears among the inputs. We then state the
$\Fq$-degeneration clause: the inequality proved here is exactly the inequality Weil's proof
of RH for curves consumes, and it is proved before, and independently of, that application ---
the short discriminant computation recovering $|N_1 - (q+1)| \le 2g\sqrt{q}$ is included. In
particular, in the function-field template the positivity engine is demonstrably not Weil
positivity in disguise: it exists prior to, and is consumed by, the RH proof, not conversely.
The historical lineage (Mattuck--Tate, Grothendieck, Hartshorne) is recorded, together with a
scoped remark on the diagonal Tate-curve surfaces $E_p \times E_p$ of Connes--Consani's
absolute geometry, where the inequality holds classically but the available correspondence
data is prime-blind.
\end{abstract}

\section*{Introduction and conventions}

\textbf{Status (read before citing).} This is a \textbf{re-derivation of standard material
for the record; zero novelty is claimed.} Every statement below is graduate-textbook
algebraic geometry. Its value is precisely the audit trail: the
Hodge-index/Castelnuovo--Severi inequality --- the positivity engine of the only
\textbf{positivity mechanism} that has ever proved an RH\footnote{``Positivity mechanism''
is the operative qualifier: RH for curves over finite fields was \emph{also} proved by the
elementary auxiliary-polynomial method of Stepanov (1969) \cite{Stepanov}, streamlined and
completed by Bombieri \cite{Bombieri} --- a proof consuming only Riemann--Roch on the
\emph{curve}: no surface, no Hodge index, no positivity engine. Deligne's Weil~I
amplification is likewise a distinct engine. The claim above concerns the Hodge-index
positivity route only.} --- is proved from Riemann--Roch on the surface, Serre duality, and
ampleness, \textbf{with no zeta-function, $L$-function, or RH input anywhere}.

\medskip

\noindent\textbf{Conventions.} $X$ is a smooth projective (geometrically irreducible)
surface over an algebraically closed field $k$ of \textbf{arbitrary characteristic}
(char~$p$ included; nothing below is characteristic-zero). Divisors are considered up to
numerical equivalence where stated; $K = K_X$ is the canonical divisor; $D \cdot E$ is the
intersection number; $h^i(D) = \dim_k H^i(X, \OX(D))$;
$\chi(D) = h^0 - h^1 + h^2$. ``$D$ numerically nontrivial'' means $D \cdot E \ne 0$ for some
divisor $E$. For a non-closed base field, all quantities below are stable under base change
to the algebraic closure, so the inequalities descend; we work over $k = \bar{k}$
throughout.

\section{What is being certified and why}
\label{sec:certified}

The function-field mechanism template for a Riemann hypothesis loads \emph{all} of its
Riemann-hypothesis content on a Hodge-index-type inequality. The known trap is a positivity
step that is Weil positivity in disguise, i.e.\ an inequality whose proof secretly consumes
the statement to be proved. The function-field template is the existence proof that the trap
is avoidable: there, the master inequality has a proof whose inputs are
\textbf{Riemann--Roch on the surface, Serre duality, and ampleness} --- nothing else. This
note re-derives that proof in full, itemizes the provenance of every input
(\S\ref{sec:ledger}), and states the $\Fq$-degeneration clause (\S\ref{sec:degeneration}):
the inequality proved here is exactly the inequality Weil's proof of RH for curves consumes,
and it is proved before and independently of that application. That is the non-circularity
claim, now on paper rather than recalled.

\paragraph{The published record, named rather than left implicit.}
None of \S\S\ref{sec:lemma}--\ref{sec:degeneration} is new. Milne's survey \cite{Milne16}
publishes them end to end and in the same order --- Riemann--Roch on a surface, Serre duality
and adjunction (\S1, p.~7); the ample-meets-effective lemma (Lemma~1.1, p.~8); the Hodge index
theorem proved from Riemann--Roch (Theorem~1.2, p.~8); the signature corollaries (1.3--1.4,
p.~9); the Castelnuovo--Severi inequality on $C_1 \times C_2$ (Theorem~1.5, p.~9); and the
deduction of $|N - q - 1| \le 2g\sqrt{q}$ (p.~10) --- closing with the observation that
\S\ref{sec:degeneration} restates: ``Note that, except for the last few lines, the proof is
purely geometric and takes place over an algebraically closed field. This is typical: study of
the Riemann hypothesis over finite fields suggests questions in algebraic geometry whose
resolution proves the hypothesis'' \cite[Aside~1.8, p.~10]{Milne16}.

Independently, and with a \emph{different} Weil-free input list, Bombieri's official Clay
problem description states the same thing about the same theorem: ``The algebraic index theorem
for surfaces is essentially due to Severi in 1906 [Sev, \S2, Teo.~I]. The proof uses the
Riemann--Roch theorem on $X$ and the finiteness of families of curves on $X$ of a given degree;
no other proof by algebraic methods is known up to now, although much later several authors
independently rediscovered Severi's argument'' \cite[p.~9]{BombieriRH}. His inputs are
Riemann--Roch plus finiteness of families of curves of given degree; this note's are
Riemann--Roch plus ampleness. Neither list contains the Weil conjectures. Note, because the
argument must not have it both ways: on Bombieri's reading Segre~1937, Bronowski~1938,
Mattuck--Tate~1958 and Grothendieck~1958 are \emph{rediscoveries of Severi's argument}, not
methodologically independent routes. What the record supports is that several published,
Weil-free derivations exist --- not that they are independent of one another.

The certification is not academic. The leading contemporary program on this problem adopts
exactly this route and states the same logical requirement in its own words: Connes and Consani
describe the ``Riemann--Roch strategy'' as trading the location of the zeros for the
non-positivity of a quadratic form $s(f, f)$, note that in the function-field case ``the most
conceptual proof was obtained by applying the Riemann--Roch formula on the square of the curve
defining the function field \ldots\ if one assumes the positivity of $s(f, f) > 0$ for some $f$,
it is the existence part of the Riemann--Roch theorem which yields a contradiction''
\cite[\S3]{CCRR}, and identify what is missing as ``the implementation of a Riemann--Roch
formula whose topological side is $\tfrac{1}{2} D \cdot D$'' \cite[\S3.1]{CCRR}.

\section{Inputs, with provenance (and what is not an input)}
\label{sec:inputs}

The proof uses exactly the following, each a standard theorem with no analytic or arithmetic
content about zeta or $L$-functions:

\begin{itemize}
\item \textbf{(IN1) Riemann--Roch on a surface.}
  $\chi(D) = \chi(\OX) + \tfrac{1}{2} D \cdot (D - K)$.
  [Hartshorne \cite{Hartshorne}, Thm.~V.1.6; valid over any algebraically closed field.]
\item \textbf{(IN2) Serre duality.} $h^2(D) = h^0(K - D)$.
  [Hartshorne \cite{Hartshorne}, III.7.6.]
\item \textbf{(IN3) Ample meets effective positively.} If $H$ is ample and $E$ is a nonzero
  effective divisor, then $E \cdot H > 0$. [Degree of a nonzero curve in a projective
  embedding is positive; equivalently the easy direction of Nakai--Moishezon.]
\item \textbf{(IN4) Ampleness absorbs any divisor.} For $H$ ample and $D$ arbitrary,
  $D + nH$ is ample for all $n \gg 0$. [Hartshorne \cite{Hartshorne}, Exercise~II.7.14(b);
  standard.]
\item \textbf{(IN5) Trivial positivity.} $h^0, h^1, h^2 \ge 0$, hence
  $h^0(D) + h^2(D) \ge \chi(D)$.
\item \textbf{(IN6) The intersection pairing on $\Num(X)$} is a well-defined, nondegenerate
  symmetric bilinear form on a finitely generated group. [Hartshorne \cite{Hartshorne}, V.1;
  nondegeneracy is the definition of numerical equivalence. Finite generation of $\Num(X)$
  is the N\'eron--Severi theorem --- proved by Picard-variety/finiteness arguments, no zeta
  input; it is consumed \emph{only} for $\rho < \infty$ in the signature statement
  \S\ref{sec:signature}, which sits off the load-bearing chain to
  \S\S\ref{sec:cs}--\ref{sec:degeneration}. This scoping is Grothendieck's own, in one
  sentence, in the paragraph immediately following his Th\'eor\`eme~1.1: ``Ce th\'eor\`eme
  semble utiliser, dans son \'enonc\'e, le th\'eor\`eme de N\'eron [5] impliquant que $E$ est
  de dimension finie. Il serait \'evident comment formuler le th\'eor\`eme 1 si on voulait
  ignorer le th\'eor\`eme de N\'eron'' \cite[p.~208]{Grothendieck} --- the finiteness input
  enters the \emph{statement}, not the \emph{proof}.]
\item \textbf{(IN7) Nakai--Moishezon criterion (hard direction).} On a surface, if
  $H^2 > 0$ and $H \cdot C > 0$ for every irreducible curve $C \subset X$, then $H$ is
  ample. [Hartshorne \cite{Hartshorne}, Thm.~V.1.10; the proof is cohomological
  (Riemann--Roch plus effectivity bootstrapping) --- no zeta content. Consumed only in
  \S\ref{sec:cs}, to certify $H = \xi_1 + \xi_2$ ample. An elementary route avoiding it
  exists --- $m(\xi_1 + \xi_2) = p_1^*(mP) + p_2^*(mP')$ is very ample for $m \gg 0$ as the
  Segre-composed embedding of very ample divisors on the factors, an (IN4)-grade fact ---
  but the criterion is what \S\ref{sec:cs}'s printed verification uses, so it is listed as
  the input it is. Ledger note: Grothendieck's own route \cite[\S2]{Grothendieck} uses neither
  (IN4) nor (IN7). His auxiliary Proposition~2.1 assumes \emph{only} $D^2 > 0$ --- no
  ampleness, no hyperplane section, no effectivity, and no sign condition on $D$ itself (he
  replaces $D$ by $-D$ if need be) --- and he obtains the single positive square directly from a
  hyperplane section $H$, since $l(H) > 1$ gives $H^2 > 0$ by his~(2.4). Adopting that route
  would shorten this ledger, but \S\ref{sec:cs} would still need $m(\xi_1 + \xi_2)$ very
  ample, so the honest accounting is that (IN7) can be \emph{replaced} by the Segre-embedding
  fact already named here, not that two inputs disappear.]
\end{itemize}

\textbf{Not inputs:} the zeta function of any curve or surface; point counts; the Weil
conjectures in any form; positivity of any explicit-formula functional; any spectral or
Hilbert-space statement; any characteristic-zero comparison. No step of the proof
(\S\S\ref{sec:lemma}--\ref{sec:cs}) mentions them. \S\ref{sec:degeneration}'s degeneration
clause consumes rationality and the functional equation of $Z(C, T)$ ---
Riemann--Roch-on-the-curve facts with no RH content --- and produces point counts as
\emph{outputs} ($N_1 = \Gamma \cdot \Delta$), exactly as itemized in the
\S\ref{sec:ledger} ledger; \S\ref{sec:baseline} mentions the explicit formula only inside a
strictly scoped fence, to state what the Tate-curve surfaces cannot host.

\section{Lemma (effectivity criterion; Hartshorne V.1.8 lineage)}
\label{sec:lemma}

\begin{lemma}
\label{lem:effectivity}
Let $H$ be ample on $X$ and let $D$ be a divisor with $D^2 > 0$ and $D \cdot H > 0$. Then
$nD$ is effective ($h^0(nD) > 0$) for all sufficiently large $n$.
\end{lemma}

\begin{proof}
By (IN1),
$\chi(nD) = \chi(\OX) + \tfrac{1}{2}\, nD \cdot (nD - K)
= \tfrac{n^2}{2} D^2 - \tfrac{n}{2}\, D \cdot K + \chi(\OX) \to \infty$
as $n \to \infty$, since $D^2 > 0$. By (IN2), $h^2(nD) = h^0(K - nD)$. Since
$D \cdot H > 0$, $(K - nD) \cdot H = K \cdot H - n (D \cdot H) < 0$ for all $n \gg 0$; but
$h^0(K - nD) > 0$ would make $K - nD$ linearly equivalent to an effective divisor $E \ge 0$
with $(K - nD) \cdot H = E \cdot H \ge 0$ (zero only for $E = 0$, positive otherwise by
(IN3)) --- impossible. So $h^0(K - nD) = 0$ for all $n \gg 0$, and by (IN5),
$h^0(nD) \ge \chi(nD) - h^2(nD) = \chi(nD) \to \infty$.
\end{proof}

\noindent\emph{(This is the exact locus where the hypothesis $D \cdot H > 0$ is consumed:
it kills the $h^2$ term. The compressed sketch ``$\chi(nD) \to \infty$ forces $\pm nD$
effective'' silently routes through this Lemma applied to an auxiliary ample divisor; the
full route is Theorem~\ref{thm:hodge}.)}

\section{Theorem (Hodge index; Hartshorne Thm.~V.1.9 lineage)}
\label{sec:hodge}

\begin{theorem}
\label{thm:hodge}
Let $H$ be ample on $X$ and let $D$ be a divisor with $D \cdot H = 0$ and $D$ numerically
nontrivial. Then $D^2 < 0$.
\end{theorem}

\begin{proof}
Suppose for contradiction $D^2 \ge 0$.

\textbf{Case 1: $D^2 > 0$.} By (IN4) choose $m \gg 0$ with $H' := D + mH$ ample. Then
$D \cdot H' = D^2 + m (D \cdot H) = D^2 > 0$, so Lemma~\ref{lem:effectivity} applies to $D$
with the ample divisor $H'$: $nD$ is effective and nonzero ($D$ is numerically nontrivial,
so $nD \not\equiv 0$) for some large $n$. By (IN3), $nD \cdot H > 0$, contradicting
$D \cdot H = 0$.

\textbf{Case 2: $D^2 = 0$.} Since $D$ is numerically nontrivial, pick $E$ with
$D \cdot E \ne 0$. Replace $E$ by $E' := (H^2) E - (E \cdot H) H$; then
$E' \cdot H = (H^2)(E \cdot H) - (E \cdot H)(H^2) = 0$ and
$D \cdot E' = (H^2)(D \cdot E) - (E \cdot H)(D \cdot H) = (H^2)(D \cdot E) \ne 0$
($H^2 > 0$ by (IN3): $H$ ample $\Rightarrow$ some multiple of $H$ is effective and nonzero,
and it meets $H$ positively, so $H^2 = H \cdot H > 0$). For an integer $n$, set
$D' := nD + E'$. Then $D' \cdot H = n (D \cdot H) + E' \cdot H = 0$ and
$D'^2 = n^2 D^2 + 2n (D \cdot E') + E'^2 = 2n (D \cdot E') + E'^2$. Choosing $n$ large with
the sign of $n (D \cdot E')$ positive gives $D'^2 > 0$ with $D' \cdot H = 0$. If $D'$ is
numerically nontrivial this contradicts Case~1; and $D'$ numerically trivial is impossible
since $D'^2 > 0$. Contradiction either way.

Hence $D^2 < 0$.
\end{proof}

\section{Corollary (the signature statement)}
\label{sec:signature}

\begin{corollary}
\label{cor:signature}
On $\Num(X) \otimes \mathbb{R}$ (rank $\rho$), the intersection form has signature
$(1, \rho - 1)$: it is positive definite on the line $\mathbb{R} \cdot H$ and negative
definite on $H^{\perp}$.
\end{corollary}

\begin{proof}
$H^2 > 0$ gives the positive line. For the negative-definiteness of $H^{\perp}$:
Theorem~\ref{thm:hodge} gives $x^2 < 0$ for every nonzero \textbf{rational} class
$x \in H^{\perp}$ (clear denominators and apply the Theorem to the resulting integral
divisor class; numerical nontriviality is exactly nonvanishing in $\Num$). Negative
\emph{semi}definiteness on the real span follows by density and continuity of the quadratic
form. If some nonzero real $x \in H^{\perp}$ had $x^2 = 0$, then $x$ would lie in the
radical of the semidefinite form restricted to $H^{\perp}$ (an isotropic vector of a
semidefinite form is radical); the radical is a rational subspace (the form and $H^{\perp}$
are defined over $\mathbb{Q}$), so it would contain a nonzero rational class
$y \in H^{\perp}$ with $y \cdot z = 0$ for all $z \in H^{\perp}$ and $y \cdot H = 0$ ---
i.e.\ $y$ numerically trivial (write any $w$ as $\lambda H + w^{\perp}$), contradicting
$y \ne 0$ in $\Num(X) \otimes \mathbb{Q}$.
\end{proof}

\section{Corollary (Castelnuovo--Severi on \texorpdfstring{$C \times C'$}{C x C'})}
\label{sec:cs}

Let $C, C'$ be smooth projective curves over $k$, $S = C \times C'$,
$\xi_1 = \{\mathrm{pt}\} \times C'$, $\xi_2 = C \times \{\mathrm{pt}\}$. Then
$\xi_1^2 = \xi_2^2 = 0$, $\xi_1 \cdot \xi_2 = 1$, and $H := \xi_1 + \xi_2$ is ample ((IN7)
Nakai--Moishezon: $H^2 = 2 > 0$, and every irreducible curve on $S$ meets $\xi_1$ or
$\xi_2$ positively and neither negatively). For a divisor $D$ on $S$ write the
\textbf{bidegree} $(d_1, d_2) := (D \cdot \xi_1, D \cdot \xi_2)$ (for the graph $\Gamma_f$
of a morphism $f\colon C \to C'$, $(d_1, d_2) = (1, \deg f)$).

\begin{corollary}[Castelnuovo--Severi inequality]
\label{cor:cs}
$D^2 \le 2 d_1 d_2$, with equality only if $D$ is numerically equivalent to
$d_2 \xi_1 + d_1 \xi_2$.
\end{corollary}

\begin{proof}
Set $D' := D - d_1 \xi_2 - d_2 \xi_1$. Then $D' \cdot \xi_1 = d_1 - d_1 = 0$ and
$D' \cdot \xi_2 = d_2 - d_2 = 0$, hence $D' \cdot H = 0$. By Theorem~\ref{thm:hodge} (and
its trivial extension: $D'^2 \le 0$, with $< 0$ unless $D' \equiv 0$),
\[
0 \ge D'^2
= D^2 - 2\bigl(d_1 (D \cdot \xi_2) + d_2 (D \cdot \xi_1)\bigr)
  + (d_1 \xi_2 + d_2 \xi_1)^2
= D^2 - 4 d_1 d_2 + 2 d_1 d_2
= D^2 - 2 d_1 d_2. \qedhere
\]
\end{proof}

The classical ``equivalence defect'' formulation and the positivity of Weil's trace form on
correspondences are algebraic reformulations of this inequality; the historical order
(\S\ref{sec:lineage}) ran through exactly this statement.

\section{The \texorpdfstring{$\mathbb{F}_q$}{Fq} degeneration clause: this is the
inequality Weil's proof consumes}
\label{sec:degeneration}

This section states, with the short computation, that the inequality of \S\ref{sec:cs} ---
proved above from (IN1)--(IN6) only --- is the entire positivity input of Weil's proof of RH
for curves over finite fields. \textbf{That is the point of this note:} the same inequality
that an arithmetic analogue of the mechanism would have to produce is, in its home
territory, proved with no RH input, so demanding a Riemann--Roch + ampleness derivation
(rather than a posited positivity) is a coherent, historically grounded non-circularity
contract.

Let $C$ be a smooth projective geometrically irreducible curve of genus $g$ over $\Fq$,
$S = C \times C$ (base-changed to $\overline{\mathbb{F}}_q$ for the intersection theory),
$F\colon C \to C$ the $q$-power Frobenius, $\Gamma = \Gamma_F$ its graph, $\Delta$ the
diagonal, $N_1 = \#C(\Fq)$. The three standard geometric facts consumed --- each an
intersection-theoretic computation with no zeta input:

\begin{itemize}
\item \textbf{(W1)} $\Gamma$ has bidegree $(1, q)$; $\Delta$ has bidegree $(1, 1)$.
\item \textbf{(W2)} $\Delta^2 = 2 - 2g$ (self-intersection $=$ degree of the normal bundle
  $N_{\Delta} \cong T_C$) and $\Gamma^2 = q(2 - 2g)$: via the embedding
  $(\mathrm{id}, F)\colon C \cong \Gamma \subset S$ one has
  $T_S|_{\Gamma} \cong T_C \oplus F^* T_C$, so
  $\deg N_{\Gamma/S} = \deg(T_S|_{\Gamma}) - \deg T_{\Gamma} = \deg F^* T_C
  = q \deg T_C = q(2 - 2g)$.
\item \textbf{(W3)} $N_1 = \Gamma \cdot \Delta$: $\Gamma$ and $\Delta$ meet transversally
  because $dF = 0$, so $1$ is never an eigenvalue of the tangent map; the intersection
  points are exactly the $\Fq$-rational points.
\end{itemize}

Apply \S\ref{sec:cs} to $D = r\Delta + s\Gamma$ for integers $r, s$: bidegree
$(r + s,\, r + sq)$, and with $t := \Gamma \cdot \Delta$,
\[
D^2 = r^2 (2 - 2g) + 2rst + s^2 q (2 - 2g) \le 2 (r + s)(r + sq).
\]
Expanding and simplifying: $g r^2 - rs\,(t - 1 - q) + g q s^2 \ge 0$ for all integers
$r, s$, hence (by homogeneity and density) for all real $r, s$; the discriminant condition
gives
\[
\bigl(t - (q + 1)\bigr)^2 \le 4 g^2 q,
\qquad \text{i.e.} \qquad
\boxed{\,|N_1 - (q + 1)| \le 2g \sqrt{q}\,.}
\]

Replacing $F$ by $F^k$ gives $|N_k - (q^k + 1)| \le 2g\, q^{k/2}$ for every $k \ge 1$; by
the standard power-sum/functional-equation bookkeeping (rationality of the zeta function of
$C$ and the $\alpha \mapsto q/\alpha$ pairing --- both consequences of Riemann--Roch on the
\emph{curve}, again with no RH input) this is equivalent to $|\alpha_i| = \sqrt{q}$ for
every Frobenius eigenvalue, i.e.\ RH for $C$. (The $\Longrightarrow$ direction is the
generating-function no-cancellation fact: the log-derivative generating function has residue
$-m_{\alpha}$ at $T = 1/\alpha$, so poles at distinct points of equal modulus cannot cancel;
the uniform bound then forces radius of convergence $\ge q^{-1/2}$, i.e.\
$\max_i |\alpha_i| \le \sqrt{q}$, and the $\alpha \mapsto q/\alpha$ pairing pins equality.)
Every input in this chain is listed in \S\ref{sec:ledger}; none is a zeta positivity.

\medskip

\noindent\textbf{Consistency note.} The bidegrees $(1, q^k)$ of $\Gamma_{F^k}$ are only the
\emph{main} terms $q^k + 1$ of $N_k$; the zero side $-\sum_i \alpha_i^k$ lives exclusively
in the $\Delta$-intersection (W3). The positivity engine acts on the trace side through
\S\ref{sec:cs}'s inequality --- degrees alone never see a zero. (Numerically spot-verified
on $y^2 = x^3 + x + 1$ over $\mathbb{F}_7$: bidegrees $(1, 7)$, $(1, 49)$, $(1, 343)$
against counts $5$, $55$, $380$.)

\section{Non-circularity audit (step-by-step input ledger)}
\label{sec:ledger}

\begin{center}
\begin{tabular}{@{}>{\raggedright\arraybackslash}p{3.4cm}>{\raggedright\arraybackslash}p{4.3cm}>{\raggedright\arraybackslash}p{3.9cm}>{\raggedright\arraybackslash}p{2.5cm}@{}}
\toprule
\textbf{Step} & \textbf{Uses} & \textbf{Provenance} & \textbf{Zeta/RH content} \\
\midrule
Lemma \S\ref{sec:lemma}, growth of $\chi(nD)$ & (IN1) RR on surfaces &
  Hartshorne V.1.6 & none \\
Lemma \S\ref{sec:lemma}, killing $h^2$ & (IN2) Serre duality $+$ (IN3) &
  Hartshorne III.7.6; projective degree & none \\
Theorem \S\ref{sec:hodge}, Case 1 & Lemma $+$ (IN4) $D + mH$ ample $+$ (IN3) &
  Hartshorne Ex.~II.7.14(b) & none \\
Theorem \S\ref{sec:hodge}, Case 2 & (IN6) linear algebra in $\Num(X)$ &
  Hartshorne V.1 & none \\
Corollary \S\ref{sec:signature} & Theorem $+$ $\mathbb{Q}$-rationality of $\Num$ &
  linear algebra & none \\
Corollary \S\ref{sec:cs} (CS) & Theorem on $C \times C'$ $+$ (IN7) ampleness of
  $H = \xi_1 + \xi_2$ & (IN7) Nakai--Moishezon, Hartshorne Thm.~V.1.10 (cohomological
  proof) & none \\
\S\ref{sec:degeneration} (W1)--(W3) & bidegree/normal-bundle/transversality computations &
  standard; $dF = 0$ & none \\
\S\ref{sec:degeneration}, discriminant step & \S\ref{sec:cs} applied to
  $r\Delta + s\Gamma$ & quadratic forms & none \\
\S\ref{sec:degeneration}, closing equivalence & rationality $+$ FE of $Z(C, T)$ from RR on
  the curve & Weil rationality (RR-based) & none (no RH input; RH is the \emph{output}) \\
\bottomrule
\end{tabular}
\end{center}

\medskip

\noindent\textbf{Statement of the non-circularity claim.} The
Hodge-index/Castelnuovo--Severi inequality on a surface, in any characteristic, \textbf{has a
proof} from Riemann--Roch on the surface, Serre duality, and ampleness alone --- \emph{a}
proof, not \emph{the} proof: Bombieri's Clay description records a different Weil-free route
through Riemann--Roch plus finiteness of families of curves of given degree, attributed to
Severi~1906 \cite[p.~9]{BombieriRH}, and Segre~1937 and Bronowski~1938 are two more, which
Grothendieck credits by name in the statement of his theorem --- ``Th\'eor\`eme~1.1
(Hodge--Segre--Bronowski)'' \cite[p.~208]{Grothendieck} --- alongside Hodge's analytic proof of
the same year. No zeta function, no
$L$-function, no point count, no explicit-formula positivity, and no statement equivalent to
or implied by any Riemann hypothesis appears among its inputs. In particular, in the
function-field template the positivity engine is demonstrably \textbf{not Weil positivity in
disguise}: it exists prior to, and is consumed by, the RH proof --- not conversely.

\medskip

\noindent\textbf{Scope of the claim.} This certifies non-circularity of the
\emph{function-field template's} engine. It does \emph{not} assert that any arithmetic
substrate exists on which the analogous inequality can be derived, and it confers no credit
by itself on any substrate proposal --- it is the gate such proposals must pass through, in
the form: an index inequality claimed on a substrate must be \emph{derived} from
Riemann--Roch plus effectivity/theta input on that substrate and must degenerate to the
present statement over $\Fq$.

\medskip

\noindent\textbf{Why the arithmetic Hodge index theorem is not the missing piece.} The one-line
objection to the paragraph above is that an arithmetic analogue of the Hodge index theorem is
already a theorem, so the substrate problem must be solved. It is not, and the reason is a
matter of codimension. The arithmetic Hodge index theorem is a statement about
\textbf{codimension-one} arithmetic cycles: on a regular arithmetic variety
$X \to \operatorname{Spec} \mathcal{O}_K$ of relative dimension $n$, if $\overline{M}$ is a
Hermitian line bundle whose generic-fiber class is primitive against an ample $\overline{L}$
($M_K \cdot L_K^{n-1} = 0$), then $\overline{M}^2 \cdot \overline{L}^{n-1} \le 0$, with
equality exactly when $\overline{M}$ is pulled back from the base --- Faltings \cite{Faltings}
and Hriljac \cite{Hriljac} for arithmetic surfaces, Moriwaki \cite{Moriwaki} for all $n$, and
in the adelic-line-bundle setting Yuan--Zhang \cite{YuanZhang}. Three things follow.

\begin{enumerate}
\item It is an inequality of one sign, in codimension one --- Moriwaki's Theorem~A(2) is
  $\widehat{\deg}(x \cdot L^{d-1}(x)) < 0$, the $p = 1$ instance of the Gillet--Soul\'e sign
  $(-1)^p \widehat{\deg}(x \cdot L^{d+1-2p}(x)) > 0$ \cite{GilletSoule}. It is the Arakelov
  shadow of the classical index theorem \emph{for divisors}, not a positivity statement about a
  pairing in the middle dimension.
\item Everything above codimension one --- exactly where a Weil-positivity argument on a
  surface-like arithmetic object would have to live --- is open: ``the high-codimensional case of
  the Gillet--Soul\'e conjecture is still wide open'' \cite[p.~2]{YuanZhang}, and K\"unnemann's
  survey confirms only $p = 0$, arithmetic surfaces, and Moriwaki's codimension-one part
  \cite[p.~115]{Kunnemann}.
\item Faltings--Hriljac \emph{is} the arithmetic-surface case --- ``when $X$ is an arithmetic
  surface \ldots\ the conjecture is a consequence of the Hodge index theorem of Faltings and
  Hriljac'' \cite[p.~115]{Kunnemann} --- so it cannot supply more on a surface than the
  codimension-one inequality already gives.
\end{enumerate}

\noindent The gate therefore stands as written: what a substrate must supply is not a
codimension-one negativity statement, which exists, but a middle-dimensional positivity calculus
on a doubled object, which does not.

\medskip

\noindent\textbf{The reading this section must answer, and its answer.} Deninger reads the
N\'eron--Tate positivity in the opposite direction --- as \emph{support} for a positivity route
to RH: with the height pairing on $\mathrm{CH}^1(E)_0 = E(\mathbb{Q}) \otimes \mathbb{Q}$
negative definite by N\'eron and Tate, ``the restriction of the Hodge inner product to the
$1$-eigenspace of $\theta$ would be identified with the negative of the (negative-definite)
N\'eron--Tate height pairing'', and ``hence our remark may be interpreted as support for the
Hilbert--Polya strategy of proving the Riemann hypotheses via positivity''
\cite[\S1]{DeningerHP}. Two features of his own wording contain the answer. His formalism is
\emph{conjectural} throughout --- ``would imply'', ``would also help'' --- so the identification
is an expectation about a cohomology theory that does not yet exist; and the support runs from a
known positivity in codimension one \emph{back} to the plausibility of that formalism, not
forward from any substrate to the inequality certified here. Read that way, his remark and this
note agree: the positivity engine is real, RH-free and upstream --- and still waiting for a
substrate to run on.

\medskip

\noindent\textbf{Two published precedents for this bookkeeping.} Kleiman's 1968 expos\'e keeps
the same ledger explicitly, separating what is proved outright from what is conditional: ``By the
Lefschetz fixed-point formula, $\chi = (\Delta^2)$ \ldots\ Hence, the functional equation,
derived with the aid of a Weil cohomology, is independent of cohomology. Furthermore, this
derivation does not depend on any unproved conjectures'' \cite[Rmk.~4.5, p.~385]{Kleiman}, while
his Theorem~4.7 is the template in which Hodge-type positivity sits logically upstream of RH. And
Ito, Ito and Koshikawa raise the identical dependency question about their own proof of a
Hodge-type positivity statement --- and answer it in the negative, in print: ``Let us recall again
that if $X$ is defined over a finite field, the Hodge standard conjecture for $X^2$ and the
Lefschetz standard conjecture for $X$ imply the Weil conjecture for $X$, and Theorem~1.2 justifies
the assumption. However, we did not attempt to avoid the Weil conjecture in our proof of
Theorem~1.2'' \cite[Rmk.~1.3]{IIK}. That is the stronger precedent, because it establishes both
that the question is a recognized one \emph{and} that answering it affirmatively is not
automatic. Ancona's Lemma~7.10 is the matching case where the proof \emph{does} consume the Weil
conjectures --- its first step invokes the functional-equation relation
$\overline{\alpha} = q/\alpha$ on Frobenius eigenvalues --- entirely legitimately, since Deligne
proved them, and that is precisely the contrast drawn here \cite[Lem.~7.10]{Ancona}.

\section{Lineage}
\label{sec:lineage}

\begin{itemize}
\item \textbf{F.~Severi} (1906) \cite{Severi} --- on Bombieri's reading
  \cite[p.~9]{BombieriRH} the algebraic index theorem is ``essentially due to Severi in 1906'',
  proved from Riemann--Roch on $X$ plus finiteness of families of curves of a given degree.
  Milne's genealogy runs Severi~1903 (the bi-additive form and its non-degeneracy conjecture)
  $\to$ Castelnuovo~1906 ($\sigma(D, D) \ge 0$) $\to$ Zariski~1952 (Riemann--Roch for normal
  surfaces in characteristic $p$) \cite[pp.~12--13]{Milne16}.
\item \textbf{W.~V.~D.~Hodge} (analytic, 1937) \cite{Hodge}, \textbf{B.~Segre} (algebraic,
  1937) \cite{Segre} and \textbf{J.~Bronowski} (algebraic, 1938) \cite{Bronowski} --- the index
  theorem itself, more than two decades before Mattuck--Tate and eleven years before Weil~1948.
  Grothendieck names his own theorem ``Th\'eor\`eme~1.1 (Hodge--Segre--Bronowski)''
  \cite[p.~208]{Grothendieck} and says of the positive-square half that Segre and Bronowski
  ``en donnent des d\'emonstrations assez simples, valables en toute caract\'eristique'' --- the
  characteristic-free claim is Grothendieck's, and is cited as his. Milne records the same in his
  footnote~16: ``The Hodge index theorem was first proved by analytic methods in Hodge 1937, and
  by algebraic methods in Segre 1937 and in Bronowski 1938.'' (Bibliographic warning:
  Grothendieck's own 1958 printed bibliography carries the corruption ``G.~Bronowski \ldots\
  (1958)'' and gives Segre's start page as 167. The correct entries are those below.)
\item \textbf{A.~Mattuck, J.~Tate}, Abh.\ Math.\ Sem.\ Univ.\ Hamburg 22 (1958)
  \cite{MattuckTate} --- the Castelnuovo--Severi inequality derived from Riemann--Roch on
  the product surface.
\item \textbf{A.~Grothendieck}, ``Sur une note de Mattuck--Tate,'' J.\ reine angew.\ Math.\
  200 (1958) 208--215 \cite{Grothendieck} --- the generalization to an arbitrary surface, with
  the numbered input list of his \S2 that \S\ref{sec:inputs} above parallels: (2.1) the
  Riemann--Roch inequality (derived there from the Riemann--Roch \emph{equality} plus Serre
  duality, so it already folds in (IN1), (IN2) and (IN5)), (2.2) Serre duality, (2.3)
  $l(D + D') \ge l(D')$ if $l(D) > 0$, and (2.4) $l(D) > 1 \Rightarrow D \cdot H > 0$ for a
  hyperplane section (this note's (IN3)). His (2.3) has no counterpart in (IN1)--(IN7).
\item \textbf{E.~Kani} (1984) \cite{Kani} --- an independent route to the defect inequality in
  characteristic $p$. \textbf{I.~Vainsencher and J.~F.~Voloch} (1988) \cite{VV} --- a further
  short published treatment, sharpening Kani. \textbf{E.~Hallouin and M.~Perret} \cite{HP} ---
  the modern continuation.
\item \textbf{R.~Hartshorne}, \emph{Algebraic Geometry}, GTM 52, Ch.~V \S1
  \cite{Hartshorne} --- the textbook consolidation used above: Thm.~V.1.6 (Riemann--Roch),
  Thm.~V.1.9 (Hodge index), Thm.~V.1.10 (Nakai--Moishezon, $=$ (IN7)), Exercises
  V.1.9--V.1.10 (Castelnuovo--Severi on $C \times C'$ and Weil's bound
  $|N - (q + 1)| \le 2g\sqrt{q}$ along exactly the route of \S\ref{sec:degeneration}).
  N.B.\ Hartshorne's theorem and exercise numbering are distinct series: ``Thm.~V.1.9''
  (the Hodge index theorem) and ``Ex.~V.1.9'' (the CS exercise) are different items that
  happen to share a number. Two sub-item pins in this note --- ``Cor.~V.1.8'' for the
  effectivity criterion and the ``(b)'' of ``Ex.~II.7.14'' --- are recalled at the letter
  level only; the statements themselves are re-proved in full above, so nothing
  load-bearing depends on them.
\end{itemize}

\section{The \texorpdfstring{$E_p \times E_p$}{Ep x Ep} baseline remark}
\label{sec:baseline}

The following remark is strictly scoped: it concerns only what the classical index
inequality does and does not provide on the per-prime Tate-curve surfaces of
Connes--Consani's absolute geometry \cite{ConnesConsani}. On the \textbf{diagonal} surfaces
$E_p \times E_p$ (with $E_p = \mathbb{C} / (\mathbb{Z} \oplus \mathbb{Z} \cdot i \log p /
2\pi)$ the per-prime Tate curve of \cite{ConnesConsani}; $\End(E_p) = \mathbb{Z}$ --- the coincidence $(\log p)^2 \in 4\pi^2 \mathbb{Q}$ is refuted unconditionally via Gelfond--Schneider, Theorem 5 of the companion note \emph{Products of the per-prime Tate curves of absolute geometry carry no correspondence calculus for the Weil explicit formula}; dated amendment 2026-08-27, this line previously hedged), the N\'eron--Severi group is
generically $\mathbb{Z}\xi_1 \oplus \mathbb{Z}\xi_2 \oplus \mathbb{Z}\Delta$, and the index
inequality of \S\S\ref{sec:hodge}--\ref{sec:cs} \textbf{holds classically} --- these are
classical complex abelian surfaces and nothing in
\S\S\ref{sec:lemma}--\ref{sec:cs} is special to them. But the available correspondence data
on them is \textbf{prime-blind}: the graphs of multiplication-by-$m$ carry Lefschetz data
$\deg([m] - 1) = (m - 1)^2$ --- integer polynomials in $m$ with no $\log p$ weight --- so no
intersection number on these surfaces can host the explicit formula's transcendentally
weighted prime terms without hand-inserting the answer key. The cross-prime surfaces
$E_p \times E_q$ are dead as substrate (off-diagonal N\'eron--Severi rank 2: no diagonal, no
graphs) and nothing in this note revives them. The baseline therefore contributes exactly
one sentence: the index inequality is available on the diagonal Tate-curve surfaces for
whatever classical purpose a future argument can state --- with zero prime-facing content.

\section*{Provenance}

\begin{remark*}
This note originated as a milestone deliverable (``M1, the non-circularity gate'') of an
internal research program on positivity mechanisms for the Riemann hypothesis, written
2026-08-26; the program's internal rules required this re-derivation to exist in writing
before any substrate proposal could claim credit for a Hodge-index-type step. The argument
of \S\S\ref{sec:lemma}--\ref{sec:hodge} was independently re-derived by two internal
reviewers, and the Mattuck--Tate and Grothendieck lineage items were verified to exist with
content as recalled (page range for \cite{Grothendieck} checked against Crelle/EuDML
records; the Bombieri expos\'e/volume/pages checked against the numdam record
SB\_1972-1973\_\_15\_\_234\_0 and Springer LNM 383). A subsequent independent referee pass
(same date) returned ``pass with repairs'' --- no fatal findings, three major and six minor,
all applied --- after re-deriving every proof line by line, verifying the input ledger of
\S\ref{sec:ledger} input-by-input with no hidden zeta/$L$-function/RH content found, and
passing 13 of 13 numerical checks (including the $\mathbb{F}_7$ spot check of
\S\ref{sec:degeneration}). The original note carries a traceability table mapping each
claim to internal program records (direction files, adjudication and verdict files, and a
catalogue of known barriers); those internal references, and the referee's itemized repair
log, are omitted here as program bookkeeping. Two Hartshorne sub-item pins
(``Cor.~V.1.8''; the ``(b)'' of ``Ex.~II.7.14'') remain recalled at the letter level only,
pending a physical-copy check; both statements are re-proved in full above. Standard
material re-derived for the record; zero novelty claimed.
\end{remark*}

\begin{thebibliography}{99}

\bibitem{Hartshorne}
R.~Hartshorne,
\emph{Algebraic Geometry},
Graduate Texts in Mathematics 52, Springer-Verlag, New York, 1977.

\bibitem{MattuckTate}
A.~Mattuck and J.~Tate,
On the inequality of Castelnuovo--Severi,
\emph{Abh.\ Math.\ Sem.\ Univ.\ Hamburg} \textbf{22} (1958), 295--299.

\bibitem{Grothendieck}
A.~Grothendieck,
Sur une note de Mattuck--Tate,
\emph{J.\ reine angew.\ Math.} \textbf{200} (1958), 208--215;
doi:10.1515/crll.1958.200.208.

\bibitem{Stepanov}
S.~A.~Stepanov,
On the number of points of a hyperelliptic curve over a finite prime field,
\emph{Izv.\ Akad.\ Nauk SSSR Ser.\ Mat.} \textbf{33} (1969), 1171--1181;
English translation: \emph{Math.\ USSR-Izv.} \textbf{3} (1969), 1103--1114.

\bibitem{Bombieri}
E.~Bombieri,
Counting points on curves over finite fields (d'apr\`es S.~A.~Stepanov),
\emph{S\'eminaire Bourbaki} 1972/73, Expos\'e 430,
Springer Lecture Notes in Mathematics 383 (1974), 234--241.

\bibitem{ConnesConsani}
A.~Connes and C.~Consani,
On the Absolute Geometry of $\mathrm{Spec}\,\mathbb{Z}$ and the
Fargues--Fontaine curve,
arXiv:2606.06604 (June 2026).

\bibitem{Ancona}
G.~Ancona,
Standard conjectures for abelian fourfolds,
\emph{Invent.\ Math.} \textbf{223} (2021), no.~1, 149--212;
doi:10.1007/s00222-020-00990-7; arXiv:1806.03216.
(Lemma~7.10; numbering as in arXiv:1806.03216v3, which post-dates Springer's
online publication by four days and is very likely --- but not verifiably --- the
published text.)

\bibitem{BombieriRH}
E.~Bombieri,
\emph{Problems of the Millennium: the Riemann Hypothesis},
official problem description, Clay Mathematics Institute, 2000, 11~pp.;
\url{https://www.claymath.org/wp-content/uploads/2022/05/riemann.pdf}.
Reprinted in \emph{The Millennium Prize Problems} (J.~Carlson, A.~Jaffe and
A.~Wiles, eds.), Amer.\ Math.\ Soc.\ and Clay Mathematics Institute, 2006,
107--124.

\bibitem{Bronowski}
J.~Bronowski,
Curves whose grade is not positive in the theory of the base,
\emph{J.\ London Math.\ Soc.} \textbf{13} (1938), 86--90;
doi:10.1112/jlms/s1-13.2.86.

\bibitem{CCRR}
A.~Connes and C.~Consani,
The Riemann--Roch strategy: Complex lift of the Scaling Site,
in: \emph{Advances in Noncommutative Geometry: On the Occasion of Alain
Connes' 70th Birthday} (A.~Chamseddine, C.~Consani, N.~Higson, M.~Khalkhali,
H.~Moscovici and G.~Yu, eds.), Springer, Cham, 2019, 53--125;
doi:10.1007/978-3-030-29597-4\_2; arXiv:1805.10501.

\bibitem{DeningerHP}
C.~Deninger,
The Hilbert--Polya strategy and height pairings,
in: \emph{Casimir Force, Casimir Operators and the Riemann Hypothesis}
(G.~van Dijk and M.~Wakayama, eds.), de Gruyter, Berlin, 2010, 275--284;
doi:10.1515/9783110226133.275; arXiv:1001.1621.

\bibitem{Faltings}
G.~Faltings,
Calculus on arithmetic surfaces,
\emph{Ann.\ of Math.} (2) \textbf{119} (1984), no.~2, 387--424;
doi:10.2307/2007043.

\bibitem{GilletSoule}
H.~Gillet and C.~Soul\'e,
Arithmetic analogs of the standard conjectures,
in: \emph{Motives} (Seattle, WA, 1991), Proc.\ Sympos.\ Pure Math.\ \textbf{55},
Part~1, Amer.\ Math.\ Soc., Providence, RI, 1994, 129--140;
doi:10.1090/pspum/055.1/1265527.

\bibitem{HP}
E.~Hallouin and M.~Perret,
From Hodge Index Theorem to the number of points of curves over finite fields,
preprint, arXiv:1409.2357 (2014).

\bibitem{Hodge}
W.~V.~D.~Hodge,
The base for algebraic varieties of given dimension,
\emph{J.\ London Math.\ Soc.} \textbf{12} (1937), 58--63;
doi:10.1112/jlms/s1-12.45.58.

\bibitem{Hriljac}
P.~Hriljac,
Heights and Arakelov's intersection theory,
\emph{Amer.\ J.\ Math.} \textbf{107} (1985), no.~1, 23--38;
doi:10.2307/2374455.

\bibitem{IIK}
K.~Ito, T.~Ito and T.~Koshikawa,
The Hodge standard conjecture for self-products of K3 surfaces,
\emph{J.\ Algebraic Geom.} \textbf{34} (2025), no.~2, 299--330;
doi:10.1090/jag/840; arXiv:2206.10086.
(Remark~1.3; numbering as in arXiv:2206.10086v1.)

\bibitem{Kani}
E.~Kani,
On Castelnuovo's equivalence defect,
\emph{J.\ reine angew.\ Math.} \textbf{352} (1984), 24--70;
doi:10.1515/crll.1984.352.24.

\bibitem{Kleiman}
S.~L.~Kleiman,
Algebraic cycles and the Weil conjectures,
in: \emph{Dix expos\'es sur la cohomologie des sch\'emas}, Advanced Studies in
Pure Mathematics \textbf{3}, North-Holland, Amsterdam, and Masson \& Cie,
Paris, 1968, 359--386.

\bibitem{Kunnemann}
K.~K\"unnemann,
Some remarks on the arithmetic Hodge index conjecture,
\emph{Compositio Math.} \textbf{99} (1995), no.~2, 109--128;
\url{http://www.numdam.org/item?id=CM_1995__99_2_109_0}.

\bibitem{Milne16}
J.~S.~Milne,
The Riemann hypothesis over finite fields: from Weil to the present day,
in: \emph{The Legacy of Bernhard Riemann After One Hundred and Fifty Years},
Vol.~II (L.~Ji, F.~Oort and S.-T.~Yau, eds.), Adv.\ Lect.\ Math.\ (ALM)
\textbf{35}, part~2, International Press, Somerville, MA, and Higher Education
Press, Beijing, 2016, 487--565; reprinted in \emph{ICCM Notices} \textbf{4}
(2016), no.~2, 14--52; doi:10.4310/ICCM.2016.v4.n2.a4; arXiv:1509.00797.
(Page pins are to \S1 of the arXiv version.)

\bibitem{Moriwaki}
A.~Moriwaki,
Hodge index theorem for arithmetic cycles of codimension one,
\emph{Math.\ Res.\ Lett.} \textbf{3} (1996), no.~2, 173--183;
doi:10.4310/MRL.1996.v3.n2.a4; arXiv:alg-geom/9403011.

\bibitem{Segre}
B.~Segre,
Intorno ad un teorema di Hodge sulla teoria della base per le curve di una
superficie algebrica,
\emph{Ann.\ Mat.\ Pura Appl.} (4) \textbf{16} (1937), 157--163;
doi:10.1007/BF02414291.

\bibitem{Severi}
F.~Severi,
Sulla totalit\`a delle curve algebriche tracciate sopra una superficie
algebrica,
\emph{Math.\ Ann.} \textbf{62} (1906), 194--225;
doi:10.1007/BF01449978.

\bibitem{VV}
I.~Vainsencher and J.~F.~Voloch,
On the Castelnuovo--Severi inequality,
\emph{J.\ reine angew.\ Math.} \textbf{390} (1988), 114--116;
doi:10.1515/crll.1988.390.114.
(Not read; the zb\textsc{math} review reports a sharpened inequality improving Kani
rather than a re-proof of the bare Castelnuovo--Severi bound.)

\bibitem{YuanZhang}
X.~Yuan and S.-W.~Zhang,
The arithmetic Hodge index theorem for adelic line bundles,
\emph{Math.\ Ann.} \textbf{367} (2017), no.~3--4, 1123--1171;
doi:10.1007/s00208-016-1414-1; arXiv:1304.3538. The published title carries
no roman numeral.

\bibitem{YuanZhangII}
X.~Yuan and S.-W.~Zhang,
The arithmetic Hodge index theorem for adelic line bundles II: finitely
generated fields,
preprint, arXiv:1304.3539 (2013; v2, 2021); unpublished as of August 2026.

\end{thebibliography}

\end{document}
